%───────────────────────────────────────────────────────────────────────────────
% EQUATION-STYLES.TEX — Coloreado semántico de ecuaciones (taxonomía híbrida)
%───────────────────────────────────────────────────────────────────────────────
%
% Sistema de "syntax highlighting" para matemáticas en tres ejes ortogonales:
%
%   1. ROL SINTÁCTICO — estable entre contextos
%      Operador, función, diferencial, dominio: siempre el mismo color.
%
%   2. ESTATUS EPISTÉMICO — la dimensión que la tipografía NO codifica
%      Incógnita vs dato/parámetro vs constante universal.
%      En f(x;θ): en física x es incógnita y θ dato; en ML es al revés.
%      El color hace explícita esta inversión sin texto adicional.
%
%   3. TIPO MATEMÁTICO — refuerzo sutil de la tipografía
%      Vector, tensor/matriz: la tipografía (bold, doble barra) ya lo indica;
%      el color refuerza cuando distinguir rango es semánticamente importante.
%
% Los colores se definen en catppuccin-palette.tex bajo el prefijo eq-*.
% Este archivo define únicamente el toggle y las macros.
%
% Toggle global: \eqcolorstrue / \eqcolorsfalse
% Uso en grupo local para versión B/N: {\eqcolorsfalse \begin{equation}...}
%───────────────────────────────────────────────────────────────────────────────

% ─── Toggle global ───────────────────────────────────────────────────────────
\newif\ifeqcolors
\eqcolorstrue     % \eqcolorsfalse para versión impresión B/N

% ═════════════════════════════════════════════════════════════════════════════
% MACROS
% ═════════════════════════════════════════════════════════════════════════════
\makeatletter

% ─── Comando interno ─────────────────────────────────────────────────────────
% Aplica color solo si el toggle está activo.
\newcommand{\@eqc}[2]{%
  \ifeqcolors\textcolor{#1}{#2}\else#2\fi
}

% ── Eje 1: Rol sintáctico ────────────────────────────────────────────────────
%   Macro      Alias de color    Cubre
%   ─────      ─────────────     ─────
%   \eqop      eq-op             Operadores: ∫, ∑, ∏, ∂, ∇, ∇², ∇×, lím, Ĥ
%   \eqfn      eq-fn             Funciones/mapas: f(·), L(θ), σ(·), V(r), ReLU
%   \eqdm      eq-dm             Diferenciales/medidas: dx, dt, dW, dS, dΩ
%   \eqdom     eq-dom            Dominios/conjuntos: Ω, ∂Ω, ℝⁿ, H, S²
\newcommand{\eqop}[1]{\@eqc{eq-op}{#1}}
\newcommand{\eqfn}[1]{\@eqc{eq-fn}{#1}}
\newcommand{\eqdm}[1]{\@eqc{eq-dm}{#1}}
\newcommand{\eqdom}[1]{\@eqc{eq-dom}{#1}}

% ── Eje 2: Estatus epistémico ────────────────────────────────────────────────
%   Macro      Alias de color    Cubre
%   ─────      ─────────────     ─────
%   \equ       eq-unknown        Incógnita / buscado / aprendido
%   \eqk       eq-known          Conocido / dato / fijado
%   \eqcon     eq-const          Constante universal / estructural: ℏ, kB, π
\newcommand{\equ}[1]{\@eqc{eq-unknown}{#1}}
\newcommand{\eqk}[1]{\@eqc{eq-known}{#1}}
\newcommand{\eqcon}[1]{\@eqc{eq-const}{#1}}

% ── Eje 3: Tipo matemático ───────────────────────────────────────────────────
%   Macro      Alias de color    Cubre
%   ─────      ─────────────     ─────
%   \eqvec     eq-vec            Campos vectoriales: u, E, B, J, f
%   \eqten     eq-ten            Tensores/matrices: σᵢⱼ, gμν, W, Σ
\newcommand{\eqvec}[1]{\@eqc{eq-vec}{#1}}
\newcommand{\eqten}[1]{\@eqc{eq-ten}{#1}}

% ═════════════════════════════════════════════════════════════════════════════
% MACROS COMPUESTOS
% ═════════════════════════════════════════════════════════════════════════════

% Derivada parcial coloreada
% Uso: \eqpdv{ρ}{t}  →  ∂ρ/∂t
% Numerador: incógnita. Denominador: variable independiente (dato del problema).
\newcommand{\eqpdv}[2]{%
  \eqop{\frac{\partial \equ{#1}}{\partial \eqk{#2}}}%
}

% Gradiente coloreado
% Uso: \eqgrad{u}  →  ∇u
\newcommand{\eqgrad}[1]{%
  \eqop{\nabla}\equ{#1}%
}

% Divergencia coloreada
% Uso: \eqdvg{\mathbf{u}}  →  ∇·u
% Nombre \eqdvg para no colisionar con \div del paquete physics.
\newcommand{\eqdvg}[1]{%
  \eqop{\nabla \cdot} \equ{#1}%
}

% SDE coloreada
% Uso: \eqSDE{X}{-\theta X}{\sigma}
% Estado = incógnita, drift = función, coef. difusión = conocido.
\newcommand{\eqSDE}[3]{%
  \eqdm{d}\equ{#1} = \eqfn{#2}\,\eqdm{dt} + \eqk{#3}\,\eqdm{dW}%
}

\makeatother
